\documentclass[institutional,screen,none]{../cls/ua-engineering-v6-beamer}
% !TeX encoding = UTF-8
\UASetUniversity{Universidad Austral}
\UASetFaculty{Facultad de Ingeniería}
\UASetDepartment{Departamento de Computación}
\UASetCampus{Pilar}
\UASetCareer{Ingeniería Informática}
\UASetCommission{Comisión A}
\UASetProfessor{Equipo docente}
\UASetSemester{Segundo cuatrimestre 2026}
\UASetCourse{Sistemas Distribuidos}
\UASetDate{2026}

% !TeX encoding = UTF-8
\UASetUniversity{Universidad Austral}
\UASetFaculty{Facultad de Ingeniería}
\UASetDepartment{Departamento de Computación}
\UASetCampus{Pilar}
\UASetCareer{Ingeniería Informática}
\UASetCommission{Comisión A}
\UASetProfessor{Equipo docente}
\UASetSemester{Segundo cuatrimestre 2026}
\UASetCourse{Sistemas Distribuidos}
\UASetDate{2026}


\UASetClassNumber{12}
\UASetClassTitle{Chaos Engineering y validación bajo falla}
\title{\UACourse}
\subtitle{Clase 12 --- Chaos Engineering, fault injection y testing distribuido}

\begin{document}

{
\setbeamertemplate{footline}{}
\begin{frame}[plain]
\titlepage
\end{frame}
}

\UASectionSlide{El problema de confiar sin verificar}

\begin{frame}{Diseñar no alcanza}
Ya vimos:
\begin{itemize}
\item consistencia
\item resiliencia
\item observabilidad
\end{itemize}

\pause

\begin{block}{Pregunta}
¿Cómo sabemos que el sistema realmente se comporta bien bajo falla?
\end{block}
\end{frame}

\begin{frame}{La intuición equivocada}
Muchos equipos asumen:
\begin{itemize}
\item “si el código parece correcto, va a resistir”
\item “si tenemos retries y breakers, alcanza”
\end{itemize}

\pause

\begin{alertblock}{Problema}
Los sistemas distribuidos fallan en interacción, no solo en componentes aislados.
\end{alertblock}
\end{frame}

\UASectionSlide{Chaos Engineering}

\begin{frame}{Definición}
\begin{block}{}
Chaos Engineering = disciplina experimental que estudia cómo responde un sistema ante perturbaciones controladas.
\end{block}
\end{frame}

\begin{frame}{Idea central}
No se trata de romper por romper.

\pause

Se trata de:
\begin{itemize}
\item formular hipótesis
\item introducir una perturbación
\item observar si el sistema mantiene el comportamiento esperado
\end{itemize}
\end{frame}

\begin{frame}{Ciclo experimental}
\begin{enumerate}
\item definir steady state
\item formular hipótesis
\item inyectar falla
\item medir impacto
\item aprender / corregir
\end{enumerate}
\end{frame}

\UASectionSlide{Fault Injection}

\begin{frame}{Tipos de falla inyectable}
\begin{itemize}
\item crash de servicio
\item latencia artificial
\item pérdida de paquetes
\item partición de red
\item saturación de CPU
\item errores de dependencia
\end{itemize}
\end{frame}

\begin{frame}{Qué se busca}
\begin{itemize}
\item ver si el sistema degrada con gracia
\item verificar breakers, retries, backpressure
\item detectar dependencias ocultas
\item validar hipótesis de operación
\end{itemize}
\end{frame}

\UASectionSlide{Steady State}

\begin{frame}{Steady State}
\begin{block}{}
Variable o conjunto de variables que representan comportamiento normal aceptable del sistema.
\end{block}

\pause

Ejemplos:
\begin{itemize}
\item tasa de checkout exitoso
\item error rate
\item latencia p95
\item mensajes procesados por segundo
\end{itemize}
\end{frame}

\begin{frame}{Por qué importa}
Sin steady state:
\begin{itemize}
\item no sabés qué validar
\item no sabés si el experimento “salió bien”
\item solo generás ruido
\end{itemize}
\end{frame}

\UASectionSlide{Testing distribuido}

\begin{frame}{No todo testing es igual}
\begin{itemize}
\item unit tests
\item integration tests
\item system tests
\item failure tests
\item chaos experiments
\end{itemize}
\end{frame}

\begin{frame}{Qué no alcanza}
\begin{itemize}
\item testear funciones aisladas
\item mocks perfectos
\item escenarios sin latencia ni pérdida
\end{itemize}

\pause

\begin{alertblock}{Clave}
Los bugs distribuidos viven en los intersticios: tiempo, orden, concurrencia y fallas parciales.
\end{alertblock}
\end{frame}

\UASectionSlide{Riesgo y seguridad}

\begin{frame}{Chaos seguro}
Buenas prácticas:
\begin{itemize}
\item experiments pequeños
\item blast radius controlado
\item rollback claro
\item observabilidad fuerte
\item horario / entorno adecuados
\end{itemize}
\end{frame}

\begin{frame}{Nunca hacer}
\begin{itemize}
\item experimentos sin hipótesis
\item experimentos sin métricas
\item experimentos en producción sin límites
\item experimentos sin plan de abortar
\end{itemize}
\end{frame}

\UASectionSlide{Ejemplo}

\begin{frame}{Ejemplo: checkout}
Hipótesis:
\begin{quote}
Si Payments responde con alta latencia, el sistema mantiene disponibilidad parcial y no colapsa por retries.
\end{quote}

Perturbación:
\begin{itemize}
\item +800 ms en Payments
\end{itemize}

Validación:
\begin{itemize}
\item breaker abre
\item p95 sube controladamente
\item checkout no derrumba todo el sitio
\end{itemize}
\end{frame}

\UASectionSlide{Cierre}

\begin{frame}{Idea central}
\begin{block}{}
La resiliencia no debe suponerse: debe validarse experimentalmente.
\end{block}
\end{frame}

\end{document}
